% Le Chat + Emilie

Response-adaptive randomization (RAR) in clinical trials aims to improve ethical and statistical efficiency by dynamically allocating patients to treatments based on observed outcomes. While RAR based on a target optimal allocation have been extensively studied for two-arms settings, their extension to multi-treatment experiments ($K \geq 2$) remains theoretically fragmented, with most existing methods focusing on specific algorithms or restricted target allocations.
In this paper, we introduce a unified framework for response-adaptive targeting, the $\alpha$-Rebalancing Targeting Strategies ($\alpha$RTS), which generalize the ERADE two-armed strategy of \cite{hu2009efficient}. We prove that all designs in this family share fundamental asymptotic properties: strong consistency, asymptotic normality of allocation proportions and treatment effect estimators, and asymptotic efficiency. To address sparse target regimes (where some treatments are asymptotically eliminated), we further propose $\alpha$RTS with Forced Exploration, a variant that guarantees infinite sampling for all treatments while preserving the asymptotic guarantees.
%Our analysis reveals a common asymptotic structure between dense and sparse regimes, suggesting potential for further unification.
Extensive simulations illustrate the finite-sample behavior of $\alpha$RTS variants in a 3-armed context, highlighting in particular the critical role of forced exploration in sparse settings.
%, comparing their convergence, hypothesis testing power, and robustness under various target allocations (Neyman, Tymofyeyev, RSIHR). The results highlight the critical role of forced exploration in sparse settings and the practical equivalence of asymptotically efficient designs.
%This work provides a flexible, theoretically grounded foundation for designing and analyzing response-adaptive targetting mechanisms in multi-arm trials, bridging the gap between bandit-inspired methods and classical clinical trial designs.


%% Redouane's prior version
%
% In this paper we provide a unified theoretical foundation for efficient response-adaptive targeting in multi-treatment experiments. Rather than analyzing a specific allocation algorithm, we introduce a broad class of response-adaptive targeting algorithms that imposes general conditions on the allocation probabilities. We will introduce two families of adaptive designs: rebalancing targeting strategies family, which operate after an initial burn-in phase, and a forced-exploration variant, which incorporate adaptive exploration for under-sampled treatments. We show that all designs in these families share the same fundamental asymptotic properties as efficient randomized adaptive designs in the multi-treatment setting, including consistency, asymptotic normality of allocation proportions and treatment effect estimators, and asymptotic efficiency. Moreover, the forced-exploration formulation allows the asymptotic theory to extend naturally to sparse target allocation regimes, in which some treatments are asymptotically eliminated.
